\Article{7}
Setzen wir den Wert des Integrals \(\int xx\varphi x.\mathrm{d}x\) von \(x=-\infty\) bis \(x=+\infty\) erstreckt \(=mm\), so nennen wir die Größe \(m\) den \textit{zu befürchtenden mittleren Fehler} oder einfach den \textit{mittleren Fehler} der Beobachtungen, deren unbestimmte Fehler \(x\) die relative Wahrscheinlichkeit \(\varphi x\) haben. Diese Benennung beschränken wir nicht auf unmittelbare Beobachtungen, sondern dehnen sie auch auf beliebige aus Beobachtungen abgeleitete Bestimmungen aus. Man muss sich jedoch sorgfältig davor hüten, den mittleren Fehler mit dem arithmetischen Mittel aller Fehler zu verwechseln, von dem wir in Art. 5 gesprochen haben.

\SourceLeaf{013}Es seien mehrere Beobachtungsarten oder mehrere aus Beobachtungen gewonnene Bestimmungen miteinander verglichen, denen nicht dieselbe Genauigkeit zuzuschreiben ist. Ihr relatives \textit{Gewicht} ist für uns eine Größe, die \(mm\) umgekehrt proportional ist, während die \textit{Genauigkeit} einfach als \(m\) umgekehrt proportional gilt. Damit ein Gewicht durch eine Zahl ausgedrückt werden kann, muss daher das Gewicht einer bestimmten Beobachtungsart als Einheit angenommen worden sein.

\Article{8}
Enthalten die Beobachtungsfehler einen konstanten Teil, so wird durch dessen Entfernung der mittlere Fehler vermindert, während Gewicht und Genauigkeit zunehmen. Behalten wir die Zeichen aus Art. 5 bei und bezeichnen den mittleren Fehler der berichtigten Beobachtungen mit \(m'\), so wird
\[
\begin{aligned}
m'm'&=\int x'x'\varphi'x'.\mathrm{d}x'=\int(x-k)^2\varphi x.\mathrm{d}x\\
&=\int xx\varphi x.\mathrm{d}x-2k\int x\varphi x.\mathrm{d}x+kk\int\varphi x.\mathrm{d}x\\
&=mm-2kk+kk=mm-kk.
\end{aligned}
\]
sein. Wäre aber statt des wahren konstanten Teils \(k\) eine andere Größe \(l\) von den Beobachtungen abgezogen worden, so würde das Quadrat des neuen mittleren Fehlers \(=mm-2kl+ll=m'm'+(l-k)^2\) werden.

\Article{9}
Bezeichnet \(\lambda\) einen bestimmten Koeffizienten und \(\mu\) den Wert des Integrals \(\int\varphi x.\mathrm{d}x\) von \(x=-\lambda m\) bis \(x=+\lambda m\), so ist \(\mu\) die Wahrscheinlichkeit, dass der Fehler einer Beobachtung kleiner als \(\lambda m\) ist (ohne Rücksicht auf das Vorzeichen), und ebenso \(1-\mu\) die Wahrscheinlichkeit eines Fehlers größer als \(\lambda m\). Entspricht also der Wert \(\mu=\tfrac12\) dem Wert \(\lambda m=\rho\), so kann der Fehler ebenso leicht unter \(\rho\) wie über \(\rho\) fallen, weshalb \(\rho\) zweckmäßig der \textit{wahrscheinliche} Fehler genannt werden kann. Die Beziehung zwischen den Größen \(\lambda\), \(\mu\) hängt offensichtlich von der Beschaffenheit der Funktion \(\varphi x\) ab, die meistens unbekannt ist. Es wird daher der Mühe wert sein, diese Beziehung für einige besondere Fälle näher zu betrachten.

I. Sind die Grenzen aller möglichen Fehler \(-a\) und \(+a\) und alle Fehler innerhalb dieser Grenzen gleich wahrscheinlich, so wird \(\varphi x\) zwischen den Grenzen \(x=-a\) und \(x=+a\) konstant und folglich \(=\tfrac{1}{2a}\) sein. Daher ist \(m=a\BareRoot{\tfrac13}\) und ebenso \(\mu=\lambda\BareRoot{\tfrac13}\), solange \(\lambda\) nicht größer als \(\BareRoot{3}\) ist; schließlich \(\rho=m\BareRoot{\tfrac34}=0{,}8660254m\), und die Wahrscheinlichkeit, dass ein Fehler auftritt, der nicht größer als der mittlere Fehler ist, wird \(=\BareRoot{\tfrac13}=0{,}5773503\) sein.

II. Sind wie zuvor \(-a\) und \(+a\) die Grenzen der möglichen Fehler und wird vorausgesetzt, dass die Wahrscheinlichkeit der Fehler selbst vom Fehler \(0\) an nach beiden Seiten in arithmetischer Folge abnimmt, so wird
\SourceLeaf{014}
\[
\begin{aligned}
\varphi x&=\frac{a-x}{aa},\quad\text{für Werte von }x\text{ zwischen }0\text{ und }+a\\[3pt]
\varphi x&=\frac{a+x}{aa},\quad\text{für Werte von }x\text{ zwischen }0\text{ und }-a
\end{aligned}
\]
sein. Hieraus folgt \(m=a\BareRoot{\tfrac16}\), \(\mu=\lambda\BareRoot{\tfrac23}-\tfrac16\lambda\lambda\), solange \(\lambda\) zwischen \(0\) und \(\BareRoot{6}\) liegt; schließlich \(\lambda=\BareRoot{6}-\BareRoot{(6-6\mu)}\), solange \(\mu\) zwischen \(0\) und \(1\) liegt, und folglich
\[
\rho=m(\BareRoot{6}-\BareRoot{3})=0{,}7174389m
\]
Die Wahrscheinlichkeit eines Fehlers, der den mittleren nicht übersteigt, wird in diesem Fall
\[
=\BareRoot{\tfrac32}-\tfrac16=0{,}6498299
\]
sein. III. Setzen wir die Funktion \(\varphi x\) proportional zu \(e^{-\frac{xx}{hh}}\) (was in der Natur allerdings nur näherungsweise wahr sein kann), so muss
\[
\varphi x=\frac{e^{-\frac{xx}{hh}}}{h\BareRoot{\pi}}
\]
sein, wobei \(\pi\) den halben Umfang eines Kreises mit dem Radius \(1\) bezeichnet, woraus wir ferner
\[
m=h\BareRoot{\tfrac12}
\]
ableiten (s. \textit{Disquis. generales circa seriem infinitam} etc. Art. 28). Bezeichnen wir ferner den Wert des Integrals
\[
\frac{2}{\BareRoot{\pi}}\int e^{-zz}\mathrm{d}z
\]
von \(z=0\) an mit \(\Theta z\), so wird
\[
\mu=\Theta(\lambda\BareRoot{\tfrac12})
\]
sein.
\Needspace{14\baselineskip}
Die folgende Tabelle gibt einige Werte dieser Größe an:
\begin{center}
\begin{tabular}{r|l}
\multicolumn{1}{c|}{\(\lambda\)}&\multicolumn{1}{c}{\(\mu\)}\\\hline
\(0{,}6744897\)&\(0{,}5\)\\
\(0{,}8416213\)&\(0{,}6\)\\
\(1{,}0000000\)&\(0{,}68268\mathord{\cdot}95\)\\
\(1{,}0364334\)&\(0{,}7\)\\
\(1{,}2815517\)&\(0{,}8\)\\
\(1{,}6448537\)&\(0{,}9\)\\
\(2{,}5758293\)&\(0{,}99\)\\
\(3{,}2918301\)&\(0{,}999\)\\
\(3{,}8905940\)&\(0{,}9999\)\\
\(\infty\)&\(1\)
\end{tabular}
\end{center}

\clearpage
